\documentclass{article}
\usepackage[sglblindworkshop,nonatbib]{neurips_2026}
\workshoptitle{ML for Systems 2026}

\usepackage[utf8]{inputenc}
\usepackage[T1]{fontenc}
\usepackage{hyperref}
\usepackage{url}
\usepackage{booktabs}
\usepackage{amsmath,amssymb}
\usepackage{nicefrac}
\usepackage{microtype}
\usepackage{xcolor}

\newcommand{\NEEDPROSE}[1]{\textcolor{red}{[TODO: #1]}}
\newcommand{\NEEDFIG}[1]{\textcolor{red}{[FIG PENDING: #1]}}

\title{When Do Attention Heads Change Their Minds?\\
       A Case for Reactive KV-Cache Residency}

\author{%
  Utkarsh Ranjan \\
  UC San Diego \\
  \texttt{\NEEDPROSE{email}} \\
}

\begin{document}
\maketitle

\begin{abstract}
Recent KV-cache optimizations classify attention heads once --- offline
or at prefill --- and treat the classification as fixed for the rest of
generation. We measure per-decode-step head activity across three
regimes (needle retrieval, long chain-of-thought, multi-turn recall) on
three models (1.5B--8B) and find 63--85\% of heads shift their reading
habits at least once within a single generation. Under a fair
evaluation, cheap runtime signals cannot predict these shifts well
enough to prefetch (precision 0.06--0.20 at recall 0.41--0.72); a
structurally different signal (generated-token identity) does clear a
useful bar, but only in long-CoT traces. We therefore argue for a
\emph{reactive} design --- demote cooling heads to CPU, fetch on
demand --- and validate it at page granularity over the same logged
attention against a frozen (FlexiCache-style) policy, a destructive
(ReasonAlloc-style) policy sharing reactive's exact schedule, and three
further cited reimplementations (SnapKV, uniform R-KV, ReasonAlloc).
Across five real model/regime combinations, reactive beats frozen at
every comparable matched-memory point (17/17), beats destructive
demotion at matched budget in 19/20 settings, and beats all three
baselines at matched memory in all 32 valid comparisons. Preliminary
real-system measurements (A30, Mistral-7B via a FlexiCache/vLLM shim)
confirm the direction: 1.34--2.33$\times$ throughput at 95--99\%
LongBench quality vs.\ stock FlexiCache. Broader real-system evaluation
is future work; the simulator-level results stand on their own.
\end{abstract}

\section{Introduction}

KV-cache size dominates long-context serving memory, and the dominant
response --- per-head compression or offload --- assumes each head's
role during a generation can be summarized once, offline or at prefill.
Recent measurement work on retrieval heads \cite{retrievalheadsdynamic}
casts doubt on that assumption and names dynamic per-head KV retention
as an open direction. This paper asks what that means for the cache.

\textbf{Contributions:} (1) replication of head non-stationarity on
three models and two regimes competing frozen-role methods do not
evaluate (long CoT, multi-turn); (2) an honest negative result --- cheap
per-head signals and their ensemble cannot predict wake-ups well enough
to prefetch, though a structurally different, data-driven signal
(generated-token identity) does, in CoT specifically; (3) a
characterization of wake-up burstiness --- real but not tightly
detectable; (4) a page-granularity residency simulator showing reactive
LRU + fetch-on-demand dominates frozen classification (C2) and that
\emph{reversible} offload specifically --- not just dynamism --- beats
destructive eviction (C3) and three cited published baselines, across
five real model/regime combinations; (5) preliminary real-system
confirmation on the throughput/quality direction.

\section{Setup}

\textbf{Models/regimes.} Qwen2.5-3B-Instruct, DeepSeek-R1-Distill-Qwen-1.5B,
DeepSeek-R1-Distill-Llama-8B; NIAH (needle retrieval, $\approx$5000-token
filler), long CoT (MATH-500~\cite{math500}, up to 2048 generated
tokens), and multi-turn recall. We log, at every decode step and head,
the top-$k$ attended positions and weights ($k=64$), a continuous needle
score and binary copy-paste score following
\cite{retrievalheadsdynamic}'s Eq.~1--2. 1.5B/3B measurements: RTX~2080~Ti
(11GB); 8B and the real-system follow-up: A30 (24GB).

\section{Results}

\begin{table}[h]
\centering
\small
\begin{tabular}{llc}
\toprule
Comparison (matched memory/budget) & Combos tested & Reactive wins \\
\midrule
vs.\ Frozen, at matched memory (C2) & 5 model/regime & \textbf{17 of 17} \\
vs.\ Evict, same LRU schedule, matched budget (C3) & 5 model/regime & \textbf{19 of 20} \\
vs.\ SnapKV / R-KV / ReasonAlloc, matched memory & 5 model/regime & \textbf{32 of 32} \\
\bottomrule
\end{tabular}
\caption{Headline results: how often reactive's miss rate was $\leq$ the
comparator's, at every real point where a fair (matched-memory or
matched-budget) comparison was possible. All on real data; see
Table~\ref{tab:m2b2} for the real-system throughput/quality result and
the main text for the one C3 exception. Details below.}
\label{tab:results}
\end{table}

\textbf{Heads shift during decoding.} Let $A_t$ denote the set of heads
(or top-$k$ pages) a given head attends to at decode step $t$. We track
step-to-step stability with the adjacent-step Jaccard index:
\begin{equation}
J(A_t, A_{t+1}) \;=\; \frac{|A_t \cap A_{t+1}|}{|A_t \cup A_{t+1}|}.
\label{eq:jaccard}
\end{equation}
$J=1$ means the attended set is unchanged between consecutive steps;
$J=0$ means complete turnover. Measured adjacent-step Jaccard is $0.53$
(continuous score, NIAH), within \cite{retrievalheadsdynamic}'s reported
$0.28$--$0.51$ range (Table~\ref{tab:churn}).

``Unstable'' heads (Table~\ref{tab:churn}) are identified with
FlexiCache's temporal-stability statistic, Random-Corrected Overlap,
which discounts raw top-$K$ overlap by its chance level under a
candidate pool of size $N$:
\begin{equation}
\mathrm{RCO}(A,B) \;=\; \max\!\left(0,\; \frac{\nicefrac{|A\cap B|}{K} - \nicefrac{K}{N}}{1 - \nicefrac{K}{N}}\right), \qquad K = \min(|A|,|B|).
\label{eq:rco}
\end{equation}
$\mathrm{RCO}=0$ at chance-level overlap and $\mathrm{RCO}=1$ when the
two sets coincide exactly. A head is flagged unstable when its mean RCO
over a trailing 16-step window falls below $0.5$. Under this test, the
fraction of heads with $\geq 1$ HeteroCache-style drift event (a
windowed-median overlap-vs-baseline dropping below $\tau{=}0.5$; the
baseline resets after each event) ranges $0.63$--$0.85$ across regimes
and holds at 8B ($0.79$ on CoT), where a small chronically-unstable tail
($2$--$9\%$) coexists with a large majority of occasionally-shifting
heads.

\begin{figure}[h]
\centering
\NEEDFIG{per-step per-head activity heatmap, CoT run, 12 most-active
heads --- visualizes the churn Table~\ref{tab:churn} quantifies}
\vspace{2.6cm}
\caption{Per-head attention-target activity over decode steps (CoT,
R1-Distill-1.5B). Each row is one head; color marks which top-$k$ pages
it attends to at each step. Bands of stable color would indicate a
frozen-role head; the visible reshuffling is what Eq.~\ref{eq:jaccard}
and Table~\ref{tab:churn} summarize numerically.}
\label{fig:churn-heatmap}
\end{figure}

\begin{table}[h]
\centering
\small
\begin{tabular}{lrrrr}
\toprule
Regime & Model & Adj.~Jaccard & Unstable frac.\ & $\geq$1 drift event \\
\midrule
NIAH & Qwen2.5-3B & 0.53 (cont.) & 0.064 & 0.63 \\
CoT & R1-Distill-1.5B & 0.82 (bin.) & 0.023 & 0.85 \\
Multi-turn & Qwen2.5-3B & 0.79 (bin.) & 0.085 & 0.76 \\
\midrule
CoT & \textbf{R1-Distill-8B} & \textbf{0.93 (bin.)} & \textbf{0.042} & \textbf{0.79} \\
\bottomrule
\end{tabular}
\caption{Head-shift metrics by regime. Adjacent Jaccard: continuous
score (NIAH), binary score (CoT/multi-turn). At 8B (bottom row), binary
Jaccard tightens as copy behavior stabilizes on math CoT, but the
page-level drift metric --- what a residency controller reacts to ---
stays high at 0.79, so the case is undiminished at scale.}
\label{tab:churn}
\end{table}

\textbf{Cheap prediction fails, honestly.} We grade four candidate
signals (page-overlap change, windowed-median drift, entropy trend,
needle-mass delta) under a \emph{causal} per-head z-score, computed from
only the trailing $w{=}32$-step window's mean $\mu$ and standard
deviation $\sigma$ (never future values, so the score is usable online):
\begin{equation}
z_t \;=\; \frac{x_t - \mu_{t-w:t}}{\sigma_{t-w:t}}.
\label{eq:zscore}
\end{equation}
This is the deployable middle ground between an oracle per-head
threshold (tunes a cutoff \emph{after} seeing labels --- not causal) and
a single global cutoff (too strict, since heads churn at different
baseline rates). Under Eq.~\ref{eq:zscore}, the best signal achieves
precision $0.08$--$0.20$ at
recall $0.41$--$0.70$: over $80\%$ of alarms are false while
$30$--$60\%$ of real wake-ups are missed. An ensemble (vote $\geq 2$ of
4) never beats the single best signal. A structurally different
candidate --- discovering, via a Bonferroni-corrected two-proportion
$z$-test on a discovery half of the data and evaluating on a held-out
half (no hand-picked marker vocabulary), whether the just-generated
token predicts a wake burst --- clears a real bar: held-out F1
$0.53$--$0.62$ on CoT, $4$--$5\times$ every attention-derived signal, but
zero markers survive on NIAH/multi-turn. Wake-ups are also temporally
bursty (mean concentration $z=+9.7$ vs.\ shuffle null on CoT) but not
tightly enough for a coarse trigger: a 2\%-of-steps trigger catches only
54\% of wake-ups. Both cheap proactive routes are ruled out.

\textbf{Reactive residency wins.} We simulate, at page granularity over
the same logged attention: \emph{full} (ceiling), \emph{frozen}
(FlexiCache-style, offline stable/unstable split), \emph{reactive}
(WakeKV: LRU-capped per-head budget, CPU-reservoir fetch-on-demand,
nothing destroyed), and \emph{evict} (identical LRU schedule to
reactive, same budget, but destructive demotion --- isolating
reversibility as a variable). Sweeping the per-head budget $B$ traces
each policy's own (memory, miss-rate) curve; since frozen and reactive
reach different memory at the same nominal $B$, we interpolate
reactive's curve, piecewise-linearly between its two bracketing
measured points, onto each of frozen's $17$ actually-measured memory
values (one per swept $B$, pooled across the five combos; $3$ of the
$20$ possible points are dropped where reactive's own sweep never
reached that much memory, not a loss). At each of those $17$ real,
matched memory levels, reactive's interpolated miss rate is $\leq$
frozen's: \textbf{17/17}. At matched budget (no interpolation needed --- both
share the identical LRU cap), reactive beats evict in \textbf{19/20}
settings; the one exception (8B/CoT, tightest budget) is explained by
resident-set composition drift after the first page reactive recovers
that evict destroyed, not a flaw in the comparison.

\begin{figure}[h]
\centering
\NEEDFIG{memory-vs-miss-rate Pareto curves, reactive vs.\ frozen vs.\
evict, overlaid for one NIAH and one CoT combo}
\vspace{2.6cm}
\caption{Matched-memory/matched-budget Pareto fronts. Reactive's curve
sits at or below frozen's at every one of the 17 matched memory points
and below evict's at 19 of 20 matched budgets (Table~\ref{tab:results}),
i.e.\ lower miss rate for the same resident footprint.}
\label{fig:pareto}
\end{figure}

\textbf{Beating published baselines.} We further implement three
faithful, cited reimplementations: SnapKV~\cite{snapkv} (one-shot
prefill selection by observation-window score, frozen through decode),
uniform R-KV~\cite{rkv} (periodic importance-ranked destructive
pruning), and ReasonAlloc~\cite{reasonalloc} (online per-head
reallocation with a starvation floor) --- each with a declared
simplification where logged data doesn't support the full method (e.g.\
no redundancy term, since raw keys aren't logged). Comparing at the same
nominal budget label is misleading here, since these three only prune
periodically and can silently use several times the memory; under the
same matched-memory correction used for C2, reactive wins \textbf{32/32}
valid comparisons. (One combo/policy pair --- Mistral-7B/NIAH vs.\
R-KV/ReasonAlloc --- is explicitly excluded, not counted: the default
128-step pruning interval never triggered on that model's short NIAH
answers, so the reported result is not a real operating point.)

\textbf{Preliminary real-system confirmation.} A FlexiCache/vLLM shim
(two config overrides: zero ``stable'' heads, rerank every step) turns
stock FlexiCache into WakeKV reactive. On Mistral-7B (LEval $8$k
$\to$ $1000$, A30), reactive beats stock throughput at every rerank
interval (Table~\ref{tab:m2b2}) while holding $95$--$99\%$ of LongBench
mean quality with no monotonic erosion. Part of the win is doing less
attention work per step (0 vs.\ 64 dense heads), not purely better
memory management: holding rerank frequency fixed at stock's native
cadence, dropping to 0 dense heads alone yields $2.33\times$ throughput
for a 5-point quality cost.

\begin{table}[h]
\centering
\small
\begin{tabular}{lrrrr}
\toprule
Config & Rerank int.\ & Gen tok/s & vs stock & LB \% \\
\midrule
Stock FlexiCache & 16 (native) & 114.4 & 1.00$\times$ & 100 \\
Reactive (WakeKV) & 1 & 153.6 & 1.34$\times$ & 97.4 \\
Reactive & 2 & 193.3 & 1.69$\times$ & 97.1 \\
Reactive & 4 & 226.3 & 1.98$\times$ & 99.1 \\
Reactive & 8 & 254.2 & 2.22$\times$ & 97.9 \\
Reactive & 16 & 266.6 & 2.33$\times$ & 95.0 \\
\bottomrule
\end{tabular}
\caption{Real-system throughput$\times$quality sweep (A30, Mistral-7B,
LEval 8k$\to$1000, 24 prompts; LongBench 4 tasks $\times$ 30 samples).
$R{=}1$ is the WakeKV design point; $R{=}16$ is the Pareto knee here
(largest throughput win, quality still $\geq 95\%$).}
\label{tab:m2b2}
\end{table}

\section{Related Work}

Static head-role methods (DuoAttention~\cite{duoattention},
RazorAttention~\cite{razorattention}, MoA~\cite{moa},
HeadKV~\cite{headkv}, FlexiCache~\cite{flexicache},
HeteroCache~\cite{heterocache}) fix each head's treatment offline or at
prefill. Adaptive allocators (Ada-KV~\cite{adakv}, CAKE~\cite{cake},
LAVa~\cite{lava}, DynamicKV~\cite{dynamickv}) are prefill-only despite
the framing. Decode-time compressors (R-KV, G-KV~\cite{gkv},
ThinKV~\cite{thinkv}, SCOPE~\cite{scope}) evict under uniform per-head
budgets. ReasonAlloc reallocates budgets during decode but destroys KV
on demotion. Offload systems (ArkVale~\cite{arkvale},
ShadowKV~\cite{shadowkv}, InfiniGen~\cite{infinigen},
HeadInfer~\cite{headinfer}, FlexiCache, HeteroCache) keep KV
recoverable but never change head priority at runtime. WakeKV closes the
intersection: dynamic \emph{and} reversible.

\section{Limitations}

Scout-scale models (1.5B--8B); five CoT runs (407 wake events) but only
one multi-turn run (73 events, coarser). The 2602.11162 binary
retrieval score is thresholded, so adjacent-step Jaccard reads higher
under it than under the continuous score --- we report both in
Table~\ref{tab:churn}. The simulator counts \emph{misses}, not
stall \emph{time}, works at page granularity approximated from top-$k$,
and compares against FlexiCache-\emph{style} policies rather than their
exact code. The real-system result is narrow (single model, single
regime, one A30) and explicitly preliminary. A real-system
reversible-vs-evict head-to-head and real-system baseline comparisons
are not yet done: wiring reactive's simulator-validated evict mode into
FlexiCache's own promote/demote path needs patching an emergent,
cross-request property of their continuous-batching internals (design
documented, not yet implemented given the risk of an unverified,
silently-wrong live-tensor patch). Broader real-system scale, workload,
and hardware coverage is future work; the honest negative prediction
result and the simulator-level Pareto/reversibility/baseline results
stand on their own.

\begin{thebibliography}{99}

\bibitem{retrievalheadsdynamic}
    \NEEDPROSE{full title and author list}, arXiv:2602.11162 (ACL 2026).
\bibitem{flexicache}
    \NEEDPROSE{full title and author list}, arXiv:2511.00868v2 (MLSys 2026).
\bibitem{heterocache}
    \NEEDPROSE{full title and author list}, arXiv:2601.13684v2 (ACL 2026 Main).
\bibitem{reasonalloc}
    Liu et al., \NEEDPROSE{full title and author list}, arXiv:2606.11164v1.
\bibitem{snapkv}
    Li et al., \NEEDPROSE{full title and author list}, arXiv:2404.14469 (NeurIPS 2024).
\bibitem{rkv}
    Cai et al., \NEEDPROSE{full title and author list}, arXiv:2505.24133.
\bibitem{duoattention}
    \NEEDPROSE{full title, author list, arXiv id, venue}.
\bibitem{razorattention}
    \NEEDPROSE{full title, author list, arXiv id, venue}.
\bibitem{moa}
    \NEEDPROSE{full title, author list, arXiv id, venue}.
\bibitem{headkv}
    \NEEDPROSE{full title, author list, arXiv id, venue}.
\bibitem{adakv}
    \NEEDPROSE{full title, author list, arXiv id, venue}.
\bibitem{cake}
    \NEEDPROSE{full title, author list, arXiv id, venue}.
\bibitem{lava}
    \NEEDPROSE{full title, author list, arXiv id, venue}.
\bibitem{dynamickv}
    \NEEDPROSE{full title, author list, arXiv id, venue}.
\bibitem{gkv}
    \NEEDPROSE{full title, author list, arXiv id, venue}.
\bibitem{thinkv}
    \NEEDPROSE{full title, author list, arXiv id, venue}.
\bibitem{scope}
    \NEEDPROSE{full title, author list, arXiv id, venue}.
\bibitem{arkvale}
    \NEEDPROSE{full title, author list, arXiv id, venue}.
\bibitem{shadowkv}
    \NEEDPROSE{full title, author list, arXiv id, venue}.
\bibitem{infinigen}
    \NEEDPROSE{full title, author list, arXiv id, venue}.
\bibitem{headinfer}
    \NEEDPROSE{full title, author list, arXiv id, venue}.
\bibitem{math500}
    \NEEDPROSE{full title, author list, arXiv id, venue}.

\end{thebibliography}

\end{document}
